\section{租赁住房不是单一“防御资产”：三只产品共享租金、续约与政策簇}

508068、508077 与 508055 均以租赁住房为核心，但收入机制、地域、价格约束、扩募历史与租约风险并不相同。三者不能因为同属“住房 REIT”而被假定为独立，亦不能把较高的期末出租率等同于可持续分派。三只 Q2 数据均为\textbf{未经审计}，且扩募或不足完整年度会破坏同比和年化比较。

\section{508068：北京保障房，稳定性与提价弹性受限并存}

508068（华夏北京保障房REIT）穿透至北京六处公共租赁住房，现金流为政府指导租金。2025 年报（经审计）披露收入 CNY104.847m、净利润 CNY45.376m、经营现金流 CNY72.049m、分派率 2.53\%、每份 NAV CNY2.9927。Q2 披露合计 4,264 套、可出租面积 241,509.97 平方米、出租率 94.07\%、收缴率 96.30\%，且披露租金标准约为 CNY35--65/平方米/月、一个租期内不调整。

\begin{exhibitbox}[表 8：508068 产品卡]
\small
\begin{tabularx}{\exhibitboxwidth}{L{3.35cm}X}
\toprule
\textbf{维度} & \textbf{核验与风险含义} \\
\midrule
Q2 经营/分派 & 收入 CNY32.292m、经营现金流 CNY24.876m、可分配金额 CNY22.430m、实际分派 CNY43.453m；管理人列示年化率 3.13\% 仅为历史披露。\\
扩募与结构 & 2025 年扩募增加资产，资产基数发生变化；两个项目公司聚合资产 CNY1.973bn、负债 CNY1.690bn，燕保宜居长期应付款 CNY836.340m。不能由此直接计算基金层杠杆。\\
市场/NAV 诊断 & 2026-08-12 最后价 CNY3.863、成交额 CNY3.198m；相对 2025-12-31 NAV 的历史诊断约 +29.08\%，相隔 224 天，非当前估值。\\
关键补证 & 政策指导租金/准入/补贴变化、扩募前后同口径现金流、项目级空置与收缴、债务期限和资本开支。\\
行动 & \qzero；不因“保障房”名称把政策稳定性误写为固定现金流。\\
\bottomrule
\end{tabularx}
\sourcenote{508068 年报、Q2、法律文件与上交所快照：\sourceid{S16--S17, S27, S28}。}
\end{exhibitbox}

\section{508077：市场化上海长租，短租期与租金下行是核心变量}

508077（华夏华润有巢租赁住房REIT）持有上海泗泾、东部经开区及马桥三处市场化长租公寓，收入还包括停车和增值服务。2025 年报（经审计）披露收入 CNY79.052m、净利润 CNY11.480m、经营现金流 CNY52.832m、分派率 1.80\%、每份 NAV CNY2.3431。Q2 披露按面积出租率 94.98\%、平均租金 CNY2.12/平方米/天，同比 -2.30\%，加权剩余租期仅 0.49 年；高收缴率 99.92\% 不能取代续租和再出租压力测试。

\begin{exhibitbox}[表 9：508077 产品卡]
\small
\begin{tabularx}{\exhibitboxwidth}{L{3.35cm}X}
\toprule
\textbf{维度} & \textbf{核验与风险含义} \\
\midrule
Q2 经营/分派 & 收入 CNY34.868m、经营现金流 CNY21.148m、可分配金额 CNY22.767m、实际分派 CNY22.539m；管理人列示年化率 3.48\% 不是预测。\\
扩募与结构 & 2025 年末扩募使可比资产基数改变；项目公司聚合资产 CNY1.630bn、负债 CNY1.700bn，两个资产披露长期应付款。必须先做结构勾稽。\\
市场/NAV 诊断 & 2026-08-12 最后价 CNY2.666、成交额 CNY15.356m；相对 2025 年末 NAV 的历史诊断约 +13.78\%，不等于当前价格结论。\\
关键补证 & 续租率、有效租金、免租优惠、各项目新签/续签节奏、扩募后同口径经营序列、债务和资本开支。\\
行动 & \qzero；租金同比下降和短 WALT 不能只用期末出租率覆盖。\\
\bottomrule
\end{tabularx}
\sourcenote{508077 年报、Q2、法律文件与上交所快照：\sourceid{S18--S19, S27, S28}。}
\end{exhibitbox}

\section{508055：上海保租房，短历史与资产内部分化要先处理}

508055（汇添富上海地产租赁住房REIT）包括江月路和虹桥两处保障性租赁住房，现金流来自租金、配套商业和停车。2025 年报是从 3 月 21 日起的部分首年，披露收入 CNY63.123m、净利润 CNY16.626m、经营现金流 CNY51.074m、分派率 2.90\%、每份 NAV CNY2.6811；这一期间不能视为完整年度基准。Q2 组合出租率 96.43\%、平均租金 CNY2.79/平方米/天、收缴率 99.32\%，但江月路出租率 97.69\% 高于虹桥的 93.16\%，需要资产级而非组合平均数判断。

\begin{exhibitbox}[表 10：508055 产品卡]
\small
\begin{tabularx}{\exhibitboxwidth}{L{3.35cm}X}
\toprule
\textbf{维度} & \textbf{核验与风险含义} \\
\midrule
Q2 经营/分派 & 收入 CNY18.248m、经营现金流 CNY13.311m、可分配金额 CNY13.899m、实际分派 CNY25.770m；管理人列示年化率 3.27\% 是历史口径。\\
结构线索 & 项目公司聚合资产 CNY1.302bn、负债 CNY1.020bn；江月路与虹桥长期应付款分别为 CNY636.000m 与 CNY257.223m。未做基金层杠杆判断。\\
市场/NAV 诊断 & 2026-08-12 最后价 CNY3.304、成交额 CNY0.829m；相对 2025 年末 NAV 的历史诊断约 +23.23\%，但历史较短且时点错配。\\
关键补证 & 完整年度经营序列、虹桥资产的租金/出租率与资本开支、长期应付款条款、项目级税后现金。\\
行动 & \qzero；不以部分首年年化或一次分派推导长期可复投现金。\\
\bottomrule
\end{tabularx}
\sourcenote{508055 年报、Q2、法律文件与上交所快照：\sourceid{S20--S21, S27, S28}。}
\end{exhibitbox}

\section{租赁住房簇的组合含义}

三只产品共同暴露于出租率、有效租金、续约、收缴、区域供给、政策和长期应付款等变量。它们必须被视为一个共同风险簇，不能在没有相关性、压力测试和管理人/赞助人关系审查的情况下分散计数。即便单只产品完成补证，REIT 总体仍要满足五个独立项目、三个管理人/赞助人组和三个运营现金流行业的激活条件。
